% GENERATED FILE. Edit canonical lessons, then run npm run generate:books.
% canonical-source-hash: fnv1a64:6ccde954a04cfbf3
% canonical-lessons: ES-C742-dieta, ES-C742-sobrepeso, ES-C742-apetito, ES-C742-ejercicio, ES-C742-masticar

\chapter{Comer bien}
\label{ch:es-comer-bien}
\hlchaptermodality{742}

\begin{chapteropening}
\textbf{By the end of this chapter:} \emph{I can talk about diet, appetite, chewing well, being overweight and exercise.}

Complete the last lesson of chapter 742: I can talk about diet, appetite, chewing well, being overweight and exercise.
\end{chapteropening}

\section[la dieta]{la dieta — the diet}

\label{lesson:ES-C742-dieta}



\begin{quote}
Before the new one: say the Spanish for insomnia, then the Spanish for the nap.
\end{quote}

\subsection*{You'll want to know: la dieta}
\textbf{la dieta} — \textquotedblleft{}the diet\textquotedblright{}.

\begin{quote}
Estoy a dieta.
\end{quote}

\begin{quote}
\emph{I am on a diet.}
\end{quote}

\emph{Estar a dieta} — to be on a diet.

\begin{cousinweb}
\emph{dieta} comes from Latin \emph{diaeta}, which was borrowed from Greek \emph{díaita} \textquotedblleft{}way of life\textquotedblright{} (EtymDB).
\end{cousinweb}

\begin{grammarlens}[title={What you've built}]
One more word for telling, explaining and describing at B1.
\end{grammarlens}

\subsection*{Your turn}
\begin{itemize}
\raggedright
  \item \emph{Say it:} \emph{la dieta}
  \item \emph{Say it:} \emph{la dieta}, once more
  \item \emph{Say it:} a sentence of your own with \emph{la dieta}
  \item \emph{Recall:} say the Spanish for insomnia, then the Spanish for the nap, then say \emph{la dieta} again
\end{itemize}

\begin{tcolorbox}[breakable,colback=teal!4,colframe=teal!35!black,title={Before you move on}]
What does la dieta mean? (\textquotedblleft{}The diet\textquotedblright{}.) Say it once more.
\end{tcolorbox}
\section[el sobrepeso]{el sobrepeso — excess weight, being overweight}

\label{lesson:ES-C742-sobrepeso}



\begin{quote}
Before the new one: say the Spanish for the nap, then the Spanish for the diet.
\end{quote}

\subsection*{You'll want to know: el sobrepeso}
\textbf{el sobrepeso} — \textquotedblleft{}excess weight, being overweight\textquotedblright{}.

\begin{quote}
El médico dice que tengo sobrepeso.
\end{quote}

\begin{quote}
\emph{The doctor says I am overweight.}
\end{quote}

\emph{Tener sobrepeso} — to be overweight.

\begin{cousinweb}
\emph{sobrepeso} is built on \emph{sobre} \textquotedblleft{}over\textquotedblright{}, from Latin \emph{super}, and \emph{peso} \textquotedblleft{}weight\textquotedblright{}, from Latin \emph{pēnsum} \textquotedblleft{}something weighed\textquotedblright{}.
\end{cousinweb}

\begin{grammarlens}[title={What you've built}]
One more word for telling, explaining and describing at B1.
\end{grammarlens}

\subsection*{Your turn}
\begin{itemize}
\raggedright
  \item \emph{Say it:} \emph{el sobrepeso}
  \item \emph{Say it:} \emph{el sobrepeso}, once more
  \item \emph{Say it:} a sentence of your own with \emph{el sobrepeso}
  \item \emph{Recall:} say the Spanish for the nap, then the Spanish for the diet, then say \emph{el sobrepeso} again
\end{itemize}

\begin{tcolorbox}[breakable,colback=teal!4,colframe=teal!35!black,title={Before you move on}]
What does el sobrepeso mean? (\textquotedblleft{}Excess weight, being overweight\textquotedblright{}.) Say it once more.
\end{tcolorbox}
\section[el apetito]{el apetito — appetite}

\label{lesson:ES-C742-apetito}



\begin{quote}
Before the new one: say the Spanish for the diet, then the Spanish for excess weight.
\end{quote}

\subsection*{You'll want to know: el apetito}
\textbf{el apetito} — \textquotedblleft{}appetite\textquotedblright{}.

\begin{quote}
No tengo apetito hoy.
\end{quote}

\begin{quote}
\emph{I have no appetite today.}
\end{quote}

\begin{cousinweb}
\emph{apetito} comes from Latin \emph{appetītus} \textquotedblleft{}desire\textquotedblright{} (EtymDB), from \emph{appetere} \textquotedblleft{}to seek after\textquotedblright{}, built on \emph{petere} \textquotedblleft{}to aim at\textquotedblright{}.
\end{cousinweb}

\begin{grammarlens}[title={What you've built}]
One more word for telling, explaining and describing at B1.
\end{grammarlens}

\subsection*{Your turn}
\begin{itemize}
\raggedright
  \item \emph{Say it:} \emph{el apetito}
  \item \emph{Say it:} \emph{el apetito}, once more
  \item \emph{Say it:} a sentence of your own with \emph{el apetito}
  \item \emph{Recall:} say the Spanish for the diet, then the Spanish for excess weight, then say \emph{el apetito} again
\end{itemize}

\begin{tcolorbox}[breakable,colback=teal!4,colframe=teal!35!black,title={Before you move on}]
What does el apetito mean? (\textquotedblleft{}Appetite\textquotedblright{}.) Say it once more.
\end{tcolorbox}
\section[el ejercicio]{el ejercicio — exercise}

\label{lesson:ES-C742-ejercicio}



\begin{quote}
Before the new one: say the Spanish for excess weight, then the Spanish for appetite.
\end{quote}

\subsection*{You'll want to know: el ejercicio}
\textbf{el ejercicio} — \textquotedblleft{}exercise\textquotedblright{}.

\begin{quote}
Hago ejercicio tres veces por semana.
\end{quote}

\begin{quote}
\emph{I exercise three times a week.}
\end{quote}

\emph{Hacer ejercicio} — to exercise. Also a school exercise.

\begin{cousinweb}
\emph{ejercicio} comes from Latin \emph{exercitium} \textquotedblleft{}exercise\textquotedblright{} (EtymDB), from \emph{exercēre} \textquotedblleft{}to keep busy, to train\textquotedblright{}.
\end{cousinweb}

\begin{grammarlens}[title={What you've built}]
One more word for telling, explaining and describing at B1.
\end{grammarlens}

\subsection*{Your turn}
\begin{itemize}
\raggedright
  \item \emph{Say it:} \emph{el ejercicio}
  \item \emph{Say it:} \emph{el ejercicio}, once more
  \item \emph{Say it:} a sentence of your own with \emph{el ejercicio}
  \item \emph{Recall:} say the Spanish for excess weight, then the Spanish for appetite, then say \emph{el ejercicio} again
\end{itemize}

\begin{tcolorbox}[breakable,colback=teal!4,colframe=teal!35!black,title={Before you move on}]
What does el ejercicio mean? (\textquotedblleft{}Exercise\textquotedblright{}.) Say it once more.
\end{tcolorbox}
\section[masticar]{masticar — to chew}

\label{lesson:ES-C742-masticar}



\begin{quote}
Before the new one: say the Spanish for appetite, then the Spanish for exercise.
\end{quote}

\subsection*{You'll want to know: masticar}
\textbf{masticar} — \textquotedblleft{}to chew\textquotedblright{}.

\begin{quote}
Hay que masticar bien la comida.
\end{quote}

\begin{quote}
\emph{You have to chew your food well.}
\end{quote}

Spelling: \emph{mastiqué}. \emph{Mascar} is the older, everyday twin.

\begin{cousinweb}
\emph{masticar} is a learned borrowing from Late Latin \emph{masticāre} \textquotedblleft{}to chew\textquotedblright{} (EtymDB); the inherited form of the same word is \emph{mascar}.
\end{cousinweb}

\begin{grammarlens}[title={What you've built}]
One more word for telling, explaining and describing at B1.
\end{grammarlens}

\subsection*{Your turn}
\begin{itemize}
\raggedright
  \item \emph{Say it:} \emph{masticar}
  \item \emph{Say it:} \emph{masticar}, once more
  \item \emph{Say it:} a sentence of your own with \emph{masticar}
  \item \emph{Recall:} say the Spanish for appetite, then the Spanish for exercise, then say \emph{masticar} again
\end{itemize}

\begin{tcolorbox}[breakable,colback=teal!4,colframe=teal!35!black,title={Before you move on}]
What does masticar mean? (\textquotedblleft{}To chew\textquotedblright{}.) Say it once more.
\end{tcolorbox}
